% answers_short.tex —— 2012 年·四川卷·理科 · 简短答案速查
% 由 tools/make_short_answers.py 自 answers.tex 的 \gksolhead/\gkanswer 提取，
% 并与 2012/2012_数学真题_四川理_答案要点.md 逐题比对后生成，勿手改（改源 MD 后重跑）。

\begin{gknotice}
\textbf{使用说明}：本节只给最终答案便于快速核对，完整推导见解析版。答案为编写者独立推导，\textbf{非官方标准答案}。
\end{gknotice}

\gkpart{选择题}
\begin{center}\begin{tabular}{ccccccccccccc}
\toprule
题号 & 1 & 2 & 3 & 4 & 5 & 6 & 7 & 8 & 9 & 10 & 11 & 12 \\
\midrule
答案 & D & B & A & B & D & C & C & B & C & A & B & D \\
\bottomrule
\end{tabular}\end{center}

\gkpart{填空题}
\begin{itemize}[leftmargin=2.6em]
  \item[\textbf{13.}] $\{a,c,d\}$。
  \item[\textbf{14.}] $90^\circ$。
  \item[\textbf{15.}] $3$。
  \item[\textbf{16.}] ①③。
\end{itemize}

\gkpart{解答题}
\begin{itemize}[leftmargin=2.6em]
  \item[\textbf{17.}] (1) $p=\dfrac15$；(2) $\xi$ 的分布列为 $P(\xi=0)=\dfrac1{1000}$，$P(\xi=1)=\dfrac{27}{1000}$，$P(\xi=2)=\dfrac{243}{1000}$，$P(\xi=3)=\dfrac{729}{1000}$；$E\xi=\dfrac{27}{10}$
  \item[\textbf{18.}] (1) $\omega=\dfrac\pi4$，值域 $\left[-2\sqrt3,\ 2\sqrt3\right]$；(2) $\dfrac{7\sqrt6}{5}$
  \item[\textbf{19.}] (1) $\dfrac{\sqrt3}4$；(2) $\dfrac{\sqrt5}5$
  \item[\textbf{20.}] (1) $a_1=a_2=0$，或 $a_1=1+\sqrt2$、$a_2=2+\sqrt2$，或 $a_1=1-\sqrt2$、$a_2=2-\sqrt2$；(2) $n=7$ 时 $T_n$ 最大，$T_7=7-\dfrac{21}2\lg2$
  \item[\textbf{21.}] (1) $x^2-\dfrac{y^2}3=1$（$x>1$）；(2) $\left(1,\ 7+4\sqrt3\right)$（$m=2$ 的退化点对应 $7$，严格解为 $(1,7)\cup(7,7+4\sqrt3)$）
  \item[\textbf{22.}] (1) $f(n)=a^n$；(2) $a$ 的最小值为 $\sqrt{17}$；(3) 前者大于后者，证明见解析
\end{itemize}
